% GENERATED FILE. Edit canonical lessons, then run npm run generate:books.
% canonical-source-hash: fnv1a64:ee89642db18d3136
% canonical-lessons: PA-C320-daktar, PA-C320-mariz, PA-C320-ilaj, PA-C320-nabaz, PA-C320-tika

\chapter{At the Clinic}
\label{ch:pa-a2-at-the-clinic}
\hlchaptermodality{320}

\begin{chapteropening}
\textbf{By the end of this chapter:} \emph{I can tell a doctor what is wrong and follow what happens.}

Complete the last lesson of chapter 320: I can tell a doctor what is wrong and follow what happens.
\end{chapteropening}

\section[\texorpdfstring{\pa{ਡਾਕਟਰ} (ḍākṭar)}{ਡਾਕਟਰ (ḍākṭar)}]{\pa{ਡਾਕਟਰ} (ḍākṭar) — a doctor}

\label{lesson:PA-C320-daktar}



\begin{quote}
Before the new one: say the Punjabi for healthy, then the Punjabi for well again after an illness.
\end{quote}

\subsection*{You'll want to know: \pa{ਡਾਕਟਰ}}
\textbf{\pa{ਡਾਕਟਰ}} — \emph{ḍākṭar} — \textquotedblleft{}a doctor\textquotedblright{}.

\begin{quote}
\pa{ਕੱਲ੍ਹ} \pa{ਮੈਂ} \pa{ਡਾਕਟਰ} \pa{ਕੋਲ} \pa{ਗਿਆ।}
\end{quote}

\begin{quote}
\emph{kall maiṁ ḍākṭar kol giā.}
\end{quote}

\begin{quote}
Yesterday I went to the doctor.
\end{quote}

A masculine or feminine noun, by the person meant.

\begin{cousinweb}
\textbf{\pa{ਡਾਕਟਰ}} (\emph{ḍākṭar}) is borrowed from English \emph{doctor}.
\end{cousinweb}

\begin{grammarlens}[title={What you've built}]
One more word for this chapter.
\end{grammarlens}

\subsection*{Your turn}
\begin{itemize}
\raggedright
  \item \emph{Say it:} \emph{ḍākṭar}
  \item \emph{Say it:} \emph{ḍākṭar}, once more
  \item \emph{Say it:} \emph{kall maiṁ ḍākṭar kol giā.} — Yesterday I went to the doctor.
  \item \emph{Recall:} say the Punjabi for healthy, then the Punjabi for well again after an illness, then say \emph{ḍākṭar} again
\end{itemize}

\begin{tcolorbox}[breakable,colback=teal!4,colframe=teal!35!black,title={Before you move on}]
What does \pa{ਡਾਕਟਰ} mean? (\textquotedblleft{}A doctor\textquotedblright{}.) Say it once more.
\end{tcolorbox}
\section[\texorpdfstring{\pa{ਮਰੀਜ਼} (marīz)}{ਮਰੀਜ਼ (marīz)}]{\pa{ਮਰੀਜ਼} (marīz) — a patient}

\label{lesson:PA-C320-mariz}



\begin{quote}
Before the new one: say the Punjabi for well again after an illness, then the Punjabi for a doctor.
\end{quote}

\subsection*{You'll want to know: \pa{ਮਰੀਜ਼}}
\textbf{\pa{ਮਰੀਜ਼}} — \emph{marīz} — \textquotedblleft{}a patient\textquotedblright{}.

\begin{quote}
\pa{ਕੱਲ੍ਹ} \pa{ਹਸਪਤਾਲ} \pa{ਵਿੱਚ} \pa{ਬਹੁਤ} \pa{ਮਰੀਜ਼} \pa{ਸਨ।}
\end{quote}

\begin{quote}
\emph{kall haspatāl vich bahut marīz san.}
\end{quote}

\begin{quote}
There were a lot of patients in the hospital yesterday.
\end{quote}

A masculine or feminine noun, by the person meant.

\begin{cousinweb}
\textbf{\pa{ਮਰੀਜ਼}} (\emph{marīz}) is a loan, usually traced through Persian to Arabic \emph{marīḍ}, \textquotedblleft{}sick\textquotedblright{}.
\end{cousinweb}

\begin{grammarlens}[title={What you've built}]
One more word for this chapter.
\end{grammarlens}

\subsection*{Your turn}
\begin{itemize}
\raggedright
  \item \emph{Say it:} \emph{marīz}
  \item \emph{Say it:} \emph{marīz}, once more
  \item \emph{Say it:} \emph{kall haspatāl vich bahut marīz san.} — There were a lot of patients in the hospital yesterday.
  \item \emph{Recall:} say the Punjabi for well again after an illness, then the Punjabi for a doctor, then say \emph{marīz} again
  \item \emph{Write it:} copy \textbf{\pa{ਮਰੀਜ਼}} once, checking the whole word
\end{itemize}

\begin{tcolorbox}[breakable,colback=teal!4,colframe=teal!35!black,title={Before you move on}]
What does \pa{ਮਰੀਜ਼} mean? (\textquotedblleft{}A patient\textquotedblright{}.) Say it once more.
\end{tcolorbox}
\section[\texorpdfstring{\pa{ਇਲਾਜ} (ilāj)}{ਇਲਾਜ (ilāj)}]{\pa{ਇਲਾਜ} (ilāj) — medical treatment}

\label{lesson:PA-C320-ilaj}



\begin{quote}
Before the new one: say the Punjabi for a doctor, then the Punjabi for a patient.
\end{quote}

\subsection*{You'll want to know: \pa{ਇਲਾਜ}}
\textbf{\pa{ਇਲਾਜ}} — \emph{ilāj} — \textquotedblleft{}medical treatment\textquotedblright{}.

\begin{quote}
\pa{ਉਸ} \pa{ਹਸਪਤਾਲ} \pa{ਵਿੱਚ} \pa{ਇਲਾਜ} \pa{ਮੁਫ਼ਤ} \pa{ਹੈ।}
\end{quote}

\begin{quote}
\emph{us haspatāl vich ilāj mufat hai.}
\end{quote}

\begin{quote}
Treatment is free at that hospital.
\end{quote}

A masculine noun.

\begin{cousinweb}
\textbf{\pa{ਇਲਾਜ}} (\emph{ilāj}) is a loan, usually traced through Persian to Arabic \emph{ilāj}, \textquotedblleft{}a remedy, treatment\textquotedblright{}.
\end{cousinweb}

\begin{grammarlens}[title={What you've built}]
One more word for this chapter.
\end{grammarlens}

\subsection*{Your turn}
\begin{itemize}
\raggedright
  \item \emph{Say it:} \emph{ilāj}
  \item \emph{Say it:} \emph{ilāj}, once more
  \item \emph{Say it:} \emph{us haspatāl vich ilāj mufat hai.} — Treatment is free at that hospital.
  \item \emph{Recall:} say the Punjabi for a doctor, then the Punjabi for a patient, then say \emph{ilāj} again
  \item \emph{Write it:} copy \textbf{\pa{ਇਲਾਜ}} once, checking the whole word
\end{itemize}

\begin{tcolorbox}[breakable,colback=teal!4,colframe=teal!35!black,title={Before you move on}]
What does \pa{ਇਲਾਜ} mean? (\textquotedblleft{}Medical treatment\textquotedblright{}.) Say it once more.
\end{tcolorbox}
\section[\texorpdfstring{\pa{ਨਬਜ਼} (nabaz)}{ਨਬਜ਼ (nabaz)}]{\pa{ਨਬਜ਼} (nabaz) — the pulse}

\label{lesson:PA-C320-nabaz}



\begin{quote}
Before the new one: say the Punjabi for a patient, then the Punjabi for medical treatment.
\end{quote}

\subsection*{You'll want to know: \pa{ਨਬਜ਼}}
\textbf{\pa{ਨਬਜ਼}} — \emph{nabaz} — \textquotedblleft{}the pulse\textquotedblright{}.

\begin{quote}
\pa{ਡਾਕਟਰ} \pa{ਹਰ} \pa{ਮਰੀਜ਼} \pa{ਦੀ} \pa{ਨਬਜ਼} \pa{ਵੇਖਦਾ} \pa{ਹੈ।}
\end{quote}

\begin{quote}
\emph{ḍākṭar har marīz dī nabaz vekhdā hai.}
\end{quote}

\begin{quote}
The doctor checks every patient's pulse.
\end{quote}

A feminine noun.

\begin{cousinweb}
\textbf{\pa{ਨਬਜ਼}} (\emph{nabaz}) is a loan, usually traced through Persian to Arabic \emph{nabḍ}, \textquotedblleft{}a pulse, a throbbing\textquotedblright{}.
\end{cousinweb}

\begin{grammarlens}[title={What you've built}]
One more word for this chapter.
\end{grammarlens}

\subsection*{Your turn}
\begin{itemize}
\raggedright
  \item \emph{Say it:} \emph{nabaz}
  \item \emph{Say it:} \emph{nabaz}, once more
  \item \emph{Say it:} \emph{ḍākṭar har marīz dī nabaz vekhdā hai.} — The doctor checks every patient's pulse.
  \item \emph{Recall:} say the Punjabi for a patient, then the Punjabi for medical treatment, then say \emph{nabaz} again
  \item \emph{Write it:} copy \textbf{\pa{ਨਬਜ਼}} once, checking the whole word
\end{itemize}

\begin{tcolorbox}[breakable,colback=teal!4,colframe=teal!35!black,title={Before you move on}]
What does \pa{ਨਬਜ਼} mean? (\textquotedblleft{}The pulse\textquotedblright{}.) Say it once more.
\end{tcolorbox}
\section[\texorpdfstring{\pa{ਟੀਕਾ} (ṭīkā)}{ਟੀਕਾ (ṭīkā)}]{\pa{ਟੀਕਾ} (ṭīkā) — an injection}

\label{lesson:PA-C320-tika}



\begin{quote}
Before the new one: say the Punjabi for medical treatment, then the Punjabi for the pulse.
\end{quote}

\subsection*{You'll want to know: \pa{ਟੀਕਾ}}
\textbf{\pa{ਟੀਕਾ}} — \emph{ṭīkā} — \textquotedblleft{}an injection\textquotedblright{}.

\begin{quote}
\pa{ਕੱਲ੍ਹ} \pa{ਹਸਪਤਾਲ} \pa{ਵਿੱਚ} \pa{ਮੈਂ} \pa{ਟੀਕਾ} \pa{ਲਿਆ।}
\end{quote}

\begin{quote}
\emph{kall haspatāl vich maiṁ ṭīkā liā.}
\end{quote}

\begin{quote}
Yesterday I had an injection at the hospital.
\end{quote}

A masculine noun.

\begin{cousinweb}
\textbf{\pa{ਟੀਕਾ}} (\emph{ṭīkā}) goes back to a reconstructed \emph{*ṭikka}, \textquotedblleft{}a mark, a spot\textquotedblright{}; CDIAL lists Punjabi \emph{ṭikkā}, \textquotedblleft{}mark\textquotedblright{}, and Hindi \emph{ṭīkā}, \textquotedblleft{}caste-mark\textquotedblright{}, under it (CDIAL 5458). How it came to mean an injection is not given there.
\end{cousinweb}

\begin{grammarlens}[title={What you've built}]
One more word for this chapter.
\end{grammarlens}

\subsection*{Your turn}
\begin{itemize}
\raggedright
  \item \emph{Say it:} \emph{ṭīkā}
  \item \emph{Say it:} \emph{ṭīkā}, once more
  \item \emph{Say it:} \emph{kall haspatāl vich maiṁ ṭīkā liā.} — Yesterday I had an injection at the hospital.
  \item \emph{Recall:} say the Punjabi for medical treatment, then the Punjabi for the pulse, then say \emph{ṭīkā} again
  \item \emph{Write it:} copy \textbf{\pa{ਟੀਕਾ}} once, checking the whole word
\end{itemize}

\begin{tcolorbox}[breakable,colback=teal!4,colframe=teal!35!black,title={Before you move on}]
What does \pa{ਟੀਕਾ} mean? (\textquotedblleft{}An injection\textquotedblright{}.) Say it once more.
\end{tcolorbox}
